\documentclass[12pt]{article}
\usepackage[utf8]{inputenc}
\usepackage[T1]{fontenc}
\usepackage{hyperref}
\title{Puragi Language Resources for K-12 Curriculum Development}
\author{Compiled by Claude Sonnet 5.5 (claude-sonnet-5-5) for Andr\'as Kornai}
\date{}
\begin{document}
\maketitle

\section{Introduction}
Puragi (also called Saga or Iwaro; ISO 639-3 \texttt{pru}) is a Papuan language of the Bird's Head Peninsula, spoken in Puragi and Saga villages, Matemani District, South Sorong Regency, Southwest Papua, Indonesia. Sources give about 700 speakers (a 1991 figure) or 1,100 (Joshua Project). It is classified in South Bird's Head; whether that family belongs to Trans--New Guinea is disputed. Very little material exists, and nothing found indicates school use or teaching materials.

\section{Primary Educational Resources}
No primers, textbooks or school programmes for Puragi were found in the searches.

\section{Linguistic and Reference Materials}
\begin{itemize}
\item Wikipedia overview of the Puragi language: \url{https://en.wikipedia.org/wiki/Puragi_language}
\item South Bird's Head languages (comparative cognate tables including Puragi): \url{https://en.wikipedia.org/wiki/South_Bird%27s_Head_languages}
\item The standard description cited is de Vries (2004), ``The Puragi language'', pp.\ 137--143 in \emph{A Short Grammar of Inanwatan} (Pacific Linguistics 560); a Leiden repository item appeared in the results but its contents were not verified: \url{https://scholarlypublications.universiteitleiden.nl/access/item%3A3463909/view}
\item ASJP source entry, Vries 2004: \url{https://asjp.clld.org/sources/8789}
\end{itemize}

\section{Religious and Community Resources}
\begin{itemize}
\item Joshua Project profile (mission-oriented; Bible translation status is listed inconsistently): \url{https://joshuaproject.net/people_groups/14486}
\end{itemize}

\section{Digital and Audio Resources}
No audio, video, corpus or software resources were found.

\section{Curriculum Development Context}
The language has a very small speaker base and almost no published material beyond a short grammatical sketch and word lists. A curriculum would need to start with community-led documentation and orthography work.

\end{document}